% Generated by scripts/build_four_wave_duration_paper_table.R; do not edit.
\begin{table}[htbp]
\centering
\small
\caption{Four-wave duration model: transition rates and recent reversals}
\label{tab:four-wave-duration-ar2}
\begin{tabular}{lrr}
\toprule
 & AR(1) & AR(2) \\
\midrule
Entry rate & 2.693 & 3.206 \\
Exit rate & 4.092 & 4.775 \\
Misclassification rate & 2.646 & 2.441 \\
Initial employment rate & 49.108 & 49.102 \\
\addlinespace
Entry after a recent exit & 6.085 & 18.460 \\
Entry after continued nonemployment & 2.547 & 2.677 \\
Exit after a recent entry & 5.860 & 29.292 \\
Exit after continued employment & 3.997 & 3.989 \\
\addlinespace
Log likelihood & -2,894,475.05 & -2,892,278.65 \\
N & 316,889 & 316,889 \\
Parameters & 33 & 35 \\
\bottomrule
\end{tabular}
\par\vspace{0.5em}
\begin{minipage}{0.98\linewidth}\footnotesize
Notes: All rates are percentages. Panel A; four waves; identical samples. Entry and exit rates are survey- and posterior-risk-weighted fitted probabilities pooled across all three quarterly transitions. History-specific rates pool only transitions 2--3 and 3--4; recent means a latent status change in the preceding interval, while continued means no such change. All four reports inform the posterior risk sets; these are not raw observed-switcher rates or real-time forecasts. Groups have different duration compositions. Both models retain the same piecewise duration hazards, duration-reporting mechanisms, and AR(1) first-transition initialization. The proposed new-nonemployment onset restriction and alternative same-employer return treatment have not been implemented. Likelihood weights are normalized to N. Standard errors and calibrated significance tests are pending; no significance stars are reported.
\end{minipage}
\end{table}
